% GENERATED FILE. Edit canonical lessons, then run npm run generate:books.
% canonical-source-hash: fnv1a64:d7f70b82b8a5c1d3
% canonical-lessons: ZH-C314-zangli, ZH-C314-shengqian, ZH-C314-anwei, ZH-C314-mudi, ZH-C314-liulei

\chapter{Saying Goodbye}
\label{ch:zh-saying-goodbye}
\hlchaptermodality{314}

\begin{chapteropening}
\textbf{By the end of this chapter:} \emph{I can talk about a funeral: remembering a life, comforting someone, and tears at the graveside.}

Complete the last lesson of chapter 314: I can talk about a funeral: remembering a life, comforting someone, and tears at the graveside.
\end{chapteropening}

\section[zànglǐ]{zànglǐ — funeral}

\label{lesson:ZH-C314-zangli}



\begin{quote}
Before the new one: say the Mandarin for to hate to leave behind, then the Mandarin for to adjust to new surroundings.
\end{quote}

\subsection*{You'll want to know: zànglǐ}
\emph{zànglǐ} — \textquotedblleft{}funeral\textquotedblright{}.

\begin{quote}
\emph{dàjiā dōu cānjiā le tā de zànglǐ} — Everyone came to his funeral.
\end{quote}

Its characters come later; for now it is a word for the ear.

\begin{cousinweb}
\emph{zànglǐ} is built from \emph{zàng} 'to bury' and \emph{lǐ} 'a rite', both inherited from Middle Chinese, \emph{tsangH} and \emph{lejX} (Guangyun, in Baxter's transcription), which glosses the first 'to hide away' and the second, quoting an older dictionary, 'what is done to serve the spirits'.
\end{cousinweb}

\begin{grammarlens}[title={What you've built}]
One more word for telling what happened.
\end{grammarlens}

\subsection*{Your turn}
\begin{itemize}
\raggedright
  \item \emph{Say it:} \emph{zànglǐ}
  \item \emph{Say it:} \emph{zànglǐ}, once more
  \item \emph{Say it:} \emph{dàjiā dōu cānjiā le tā de zànglǐ} — Everyone came to his funeral.
  \item \emph{Recall:} say the Mandarin for to hate to leave behind, then the Mandarin for to adjust to new surroundings, then say \emph{zànglǐ} again
\end{itemize}

\begin{tcolorbox}[breakable,colback=teal!4,colframe=teal!35!black,title={Before you move on}]
What does zànglǐ mean? (\textquotedblleft{}Funeral\textquotedblright{}.) Say it once more.
\end{tcolorbox}
\section[\texorpdfstring{\zh{生前} (shēngqián)}{生前 (shēngqián)}]{\zh{生前} (shēngqián) — in someone's lifetime; while they were alive}

\label{lesson:ZH-C314-shengqian}



\begin{quote}
Before the new one: say the Mandarin for to adjust to new surroundings, then the Mandarin for a funeral.
\end{quote}

\subsection*{You'll want to know: \zh{生前}}
\textbf{\zh{生前}} — \emph{shēngqián} — \textquotedblleft{}in someone's lifetime; while they were alive\textquotedblright{}.

\begin{quote}
\textbf{\zh{我还记得他生前很喜欢喝茶}} \emph{wǒ hái jìde tā shēngqián hěn xǐhuan hē chá} — I still remember how he loved his tea.
\end{quote}

\begin{cousinweb}
\textbf{\zh{生前}} (\emph{shēngqián}) is built from \zh{生} 'life' and \zh{前} 'before', both inherited from Middle Chinese, \emph{srjaeng} and \emph{dzen} (Guangyun, in Baxter's transcription): before, in life, said of someone who has died.
\end{cousinweb}

\begin{grammarlens}[title={What you've built}]
One more word for telling what happened.
\end{grammarlens}

\subsection*{Your turn}
\begin{itemize}
\raggedright
  \item \emph{Say it:} \emph{shēngqián}
  \item \emph{Say it:} \emph{shēngqián}, once more
  \item \emph{Say it:} \emph{wǒ hái jìde tā shēngqián hěn xǐhuan hē chá} — I still remember how he loved his tea.
  \item \emph{Recall:} say the Mandarin for to adjust to new surroundings, then the Mandarin for a funeral, then say \emph{shēngqián} again
\end{itemize}

\begin{tcolorbox}[breakable,colback=teal!4,colframe=teal!35!black,title={Before you move on}]
What does \zh{生前} mean? (\textquotedblleft{}In someone's lifetime; while they were alive\textquotedblright{}.) Say it once more.
\end{tcolorbox}
\section[ānwèi]{ānwèi — to comfort; to console}

\label{lesson:ZH-C314-anwei}



\begin{quote}
Before the new one: say the Mandarin for a funeral, then the Mandarin for in someone's lifetime.
\end{quote}

\subsection*{You'll want to know: ānwèi}
\emph{ānwèi} — \textquotedblleft{}to comfort; to console\textquotedblright{}.

\begin{quote}
\emph{wǒ qù ānwèi tā} — I'm going to comfort her.
\end{quote}

Its characters come later; for now it is a word for the ear.

\begin{cousinweb}
\emph{ānwèi} is built from \emph{ān} 'calm' and \emph{wèi} 'to soothe', both inherited from Middle Chinese, \emph{'an} and \emph{'jw+jH} (Guangyun, in Baxter's transcription), which glosses the second with this very word.
\end{cousinweb}

\begin{grammarlens}[title={What you've built}]
One more word for telling what happened.
\end{grammarlens}

\subsection*{Your turn}
\begin{itemize}
\raggedright
  \item \emph{Say it:} \emph{ānwèi}
  \item \emph{Say it:} \emph{ānwèi}, once more
  \item \emph{Say it:} \emph{wǒ qù ānwèi tā} — I'm going to comfort her.
  \item \emph{Recall:} say the Mandarin for a funeral, then the Mandarin for in someone's lifetime, then say \emph{ānwèi} again
\end{itemize}

\begin{tcolorbox}[breakable,colback=teal!4,colframe=teal!35!black,title={Before you move on}]
What does ānwèi mean? (\textquotedblleft{}To comfort; to console\textquotedblright{}.) Say it once more.
\end{tcolorbox}
\section[mùdì]{mùdì — cemetery; graveyard}

\label{lesson:ZH-C314-mudi}



\begin{quote}
Before the new one: say the Mandarin for in someone's lifetime, then the Mandarin for to comfort someone who is sad.
\end{quote}

\subsection*{You'll want to know: mùdì}
\emph{mùdì} — \textquotedblleft{}cemetery; graveyard\textquotedblright{}.

\begin{quote}
\emph{mùdì lí zhèr hěn yuǎn} — The cemetery is a long way from here.
\end{quote}

Its characters come later; for now it is a word for the ear.

\begin{cousinweb}
\emph{mùdì} is built from \emph{mù} 'a grave' and \zh{地} 'ground', both inherited from Middle Chinese, \emph{muH} and \emph{dijH} (Guangyun, in Baxter's transcription), which glosses the first 'a grave mound'.
\end{cousinweb}

\begin{grammarlens}[title={What you've built}]
One more word for telling what happened.
\end{grammarlens}

\subsection*{Your turn}
\begin{itemize}
\raggedright
  \item \emph{Say it:} \emph{mùdì}
  \item \emph{Say it:} \emph{mùdì}, once more
  \item \emph{Say it:} \emph{mùdì lí zhèr hěn yuǎn} — The cemetery is a long way from here.
  \item \emph{Recall:} say the Mandarin for in someone's lifetime, then the Mandarin for to comfort someone who is sad, then say \emph{mùdì} again
\end{itemize}

\begin{tcolorbox}[breakable,colback=teal!4,colframe=teal!35!black,title={Before you move on}]
What does mùdì mean? (\textquotedblleft{}Cemetery; graveyard\textquotedblright{}.) Say it once more.
\end{tcolorbox}
\section[liúlèi]{liúlèi — to shed tears}

\label{lesson:ZH-C314-liulei}



\begin{quote}
Before the new one: say the Mandarin for to comfort someone who is sad, then the Mandarin for a cemetery.
\end{quote}

\subsection*{You'll want to know: liúlèi}
\emph{liúlèi} — \textquotedblleft{}to shed tears\textquotedblright{}.

\begin{quote}
\emph{dàjiā dōu liúlèi le} — Everyone was in tears.
\end{quote}

Its characters come later; for now it is a word for the ear.

\begin{cousinweb}
\emph{liúlèi} is built from \emph{liú} 'to flow' and \emph{lèi} 'tears', both inherited from Middle Chinese, \emph{ljuw} and \emph{lwijH} (Guangyun, in Baxter's transcription), which glosses the first, quoting an older dictionary, 'water moving' and the second 'tears'.
\end{cousinweb}

\begin{grammarlens}[title={What you've built}]
One more word for telling what happened.
\end{grammarlens}

\subsection*{Your turn}
\begin{itemize}
\raggedright
  \item \emph{Say it:} \emph{liúlèi}
  \item \emph{Say it:} \emph{liúlèi}, once more
  \item \emph{Say it:} \emph{dàjiā dōu liúlèi le} — Everyone was in tears.
  \item \emph{Recall:} say the Mandarin for to comfort someone who is sad, then the Mandarin for a cemetery, then say \emph{liúlèi} again
\end{itemize}

\begin{tcolorbox}[breakable,colback=teal!4,colframe=teal!35!black,title={Before you move on}]
What does liúlèi mean? (\textquotedblleft{}To shed tears\textquotedblright{}.) Say it once more.
\end{tcolorbox}
